%=== CHAPTER THREE ===
%% Kinematics, dynamics, and the parameter structure of the payload-adaptive
%% model.  2026-11-06: describes the mainline estimator only (reduced
%% velocity--mass MHE).  The full-state estimator of the earlier campaigns
%% (eq:xaug, eq:pmhe, eq:mheresid, eq:nodwdm, eq:Afrozen) is defined in
%% Chapter 5, Section sec:model:mhe.  The controller's model (eq:pnmpc,
%% eq:predmodel) is in Chapter 4, Section 4.2.
%% Restructured 2026-11-03 into six sections.  Implementation detail that used
%% to live here (initialisation, the NMPC problem and its input baseline, the
%% actuator envelope and allocation, the residual specific force, inner-loop
%% bandwidth) now lives in Chapter 4 and keeps its original labels
%% (sec:model:seed, eq:nmpc, eq:hover, eq:envelope, sec:model:residual, ...),
%% so cross-references resolve unchanged.  Experimental numbers live in
%% Chapter 5.  The pre-restructure source is in git history.

\chapter{Modeling and Parameter Structure}
\label{ch:model}
\begin{spacing}{1.5}
\setlength{\parskip}{0.22in}

A quadrotor that picks up a parcel does not merely become heavier. It also
becomes harder to rotate: on the platform of this work, a \SI{0.30}{\kg} parcel
hanging about a quarter of a metre below the rotor plane adds fifteen percent to
the mass but more than doubles the roll and pitch inertia
(Section~\ref{sec:model:magnitudes}). These are two distinct changes to the
equations of motion, and a controller that compensates only the first loses
inner-loop bandwidth in a way that mass adaptation cannot explain.

This chapter derives the model that the estimator and the controller share and
makes explicit which of its parameters are identified online. Its structural
claim is that, for a payload grasped on the thrust axis, both changes are
governed by a \emph{single} online scalar. The composite mass $\mT$ is
identified from the translational channel; the roll and pitch inertia increment
then follows from $\mT$ by rigid-body algebra and the nominal vertical arm of
the gripper. No payload mass is supplied to either optimiser, and no attach or
release notification to the estimator. The nominal arm is a geometric prior, and
Section~\ref{sec:model:whynot} states exactly what that prior is and what it is
not.

Section~\ref{sec:model:frames} defines the system and its assumptions,
Section~\ref{sec:model:dyn} its dynamics, and Section~\ref{sec:model:payload}
the coupling between payload mass and inertia. Section~\ref{sec:model:onedof}
reduces the payload to one online degree of freedom and explains why the
estimator keeps away from the rotational channel. Section~\ref{sec:model:phys}
describes the measured input and the interface to the autopilot, and
Section~\ref{sec:model:approx} collects the modelling assumptions and their
limits. The estimator and controller models built on this chapter are collected
in Section~\ref{sec:frame:models}; their numerical form, initialisation and
constraints, the residual specific force and the inner-loop scheduling belong to
Chapter~\ref{ch:framework}; the
experimental evidence for the design choices belongs to
Chapter~\ref{ch:evaluation}.

\section{System Definition and Modelling Assumptions}\label{sec:model:frames}

Two reference frames are used throughout. The world frame $\mathcal{W}$ is
East--North--Up (ENU), so that gravity acts along $-z$. The body frame
$\mathcal{B}$ is Front--Left--Up (FLU) with its origin at the centre of mass of
the \emph{bare airframe}, on the thrust axis. This is the point to which the
autopilot's odometry and the inertial unit refer; in the simulated airframe the
rotor plane lies \SI{0.032}{\meter} above it. The origin is fixed in the body,
so once a payload is attached it is \emph{not} the centre of mass of the
composite vehicle (Section~\ref{sec:model:refpoint}). The thrust vector lies on
the body $z$ axis through the origin.

The state vector and the control input are
\begin{equation}
x = \bigl[\bm{p};\, \bm{v};\, \bm{q};\, \bm{\omega}\bigr] \in \mathbb{R}^{13},
\qquad
u = \bigl[T;\, \bm{\tau}\bigr] \in \mathbb{R}^{4},
\label{eq:state}
\end{equation}
where $\bm{p}$ and $\bm{v}$ are the position and velocity in $\mathcal{W}$,
$\bm{q}$ is the unit quaternion encoding the attitude of $\mathcal{B}$ relative
to $\mathcal{W}$ (Section~\ref{sec:model:quat}),
$\bm{\omega}$ is the body rate in $\mathcal{B}$, $T$ is the collective thrust
along the body $z$ axis, and $\bm{\tau}$ is the body torque. We abbreviate
$\ev = [0,0,1]^{\!\top}$, and write $\mB$, $\mP$ and $\mT=\mB+\mP$ for the
bare-airframe, payload and composite masses.

Three assumptions define the system.
\begin{enumerate}
\item \emph{The bare airframe is balanced about the body origin}, so that its
centre of mass and the body origin coincide when no payload is attached. This is
an assumption about the vehicle, not a consequence of the frame choice.
\item \emph{The payload is a point mass $\mP$ rigidly attached on the thrust
axis}, at the body-frame offset $\rp = [0,\,0,\,r_z]^{\!\top}$ with $r_z<0$: the
load hangs below the airframe. The gripper aligns the vehicle over the parcel
before attaching; the residual horizontal offset that alignment leaves is
treated as an approximation in Section~\ref{sec:model:approx}.
\item \emph{Rotor thrust follows the static quadratic law} of
Section~\ref{sec:model:phys}, with calibrated coefficients.
\end{enumerate}

\subsection{Attitude representation}\label{sec:model:quat}

Attitude is represented by a unit quaternion under the \emph{Hamilton} convention
with the scalar part first,
\begin{equation}
\bm{q} = [q_w,\, q_x,\, q_y,\, q_z]^{\!\top},
\qquad \|\bm{q}\| = 1,
\label{eq:quatdef}
\end{equation}
encoding the rotation from $\mathcal{B}$ to $\mathcal{W}$. The corresponding
rotation matrix is
\begin{equation}
R(\bm{q}) =
\begin{bmatrix}
q_w^2 + q_x^2 - q_y^2 - q_z^2 & 2(q_xq_y - q_wq_z) & 2(q_xq_z + q_wq_y) \\[2pt]
2(q_xq_y + q_wq_z) & q_w^2 - q_x^2 + q_y^2 - q_z^2 & 2(q_yq_z - q_wq_x) \\[2pt]
2(q_xq_z - q_wq_y) & 2(q_yq_z + q_wq_x) & q_w^2 - q_x^2 - q_y^2 + q_z^2
\end{bmatrix}.
\label{eq:Rq}
\end{equation}
Three conventions are fixed here and used throughout: the Hamilton product, the
scalar-first component order of \eqref{eq:quatdef}, and $\bm{\omega}$ resolved
in $\mathcal{B}$.

One scalar derived from \eqref{eq:Rq} is used repeatedly: the fraction of the
collective thrust that acts vertically. With roll $\phi$ and pitch $\theta$ in
the ZYX Euler convention $R = R_z(\psi)R_y(\theta)R_x(\phi)$, it is the $(3,3)$
entry of the rotation matrix,
\begin{equation}
\cos\phi\,\cos\theta \;=\; \ev^{\!\top} R(\bm{q})\, \ev \;=\; 1 - 2\bigl(q_x^2+q_y^2\bigr).
\label{eq:costheta}
\end{equation}
The implementation evaluates the right-hand side directly from the quaternion.

\section{Quadrotor Dynamics with a Rigid Payload}\label{sec:model:dyn}

\subsection{Kinematics}\label{sec:model:kin}

Translational kinematics is the definition of velocity,
\begin{equation}
\dot{\bm{p}} = \bm{v}.
\label{eq:kin}
\end{equation}
With the body rate resolved in $\mathcal{B}$, attitude kinematics is the
Hamilton product of the attitude with the pure quaternion of the body rate,
\begin{equation}
\dot{\bm{q}} = \tfrac{1}{2}\, \bm{q} \otimes [0,\, \bm{\omega}]
= \tfrac{1}{2}\, \Xi(\bm{q})\, \bm{\omega},
\qquad
\Xi(\bm{q}) =
\begin{bmatrix}
-q_x & -q_y & -q_z \\
 q_w & -q_z &  q_y \\
 q_z &  q_w & -q_x \\
-q_y &  q_x &  q_w
\end{bmatrix}.
\label{eq:quat}
\end{equation}
The body rate enters on the \emph{right} of the product because it is resolved
in the body frame; a world-frame rate would enter on the left. The
implementation uses the matrix form, which is linear in $\bm\omega$.
One property of $\Xi$ matters for prediction: $\bm{q}^{\!\top}\Xi(\bm{q})=\bm 0^{\!\top}$
for every quaternion, so $\bm q^{\!\top}\dot{\bm q}=0$ and the \emph{continuous}
kinematics preserve $\lVert\bm q\rVert$ exactly. The discretised kinematics do
not: one explicit fourth-order Runge--Kutta step reproduces the exact rotation
only to fourth order in $\lVert\bm\omega\rVert\Delta t$, and the norm drifts
by a term of sixth order. The controller integrates \eqref{eq:quat} across its
horizon without renormalising and evaluates \eqref{eq:Rq} on the unnormalised
prediction, so a drift would scale the predicted thrust direction by
$\lVert\bm q\rVert^2$. At the controller's step of \SI{0.05}{\second} the drift
over the twenty-step horizon is $3\times10^{-11}$ at \SI{1}{\radian\per\second},
$2\times10^{-8}$ at \SI{3}{\radian\per\second} and $9\times10^{-6}$ at
\SI{8}{\radian\per\second}. Integrated along the measured body rates of all $96$
flights of Section~\ref{sec:eval:reduced} (largest component
\SI{1.6}{\radian\per\second}), the largest drift over any horizon was
$3\times10^{-10}$. Each solve starts from the measured unit quaternion, so the
drift does not accumulate between solves, and renormalisation was omitted as
numerically irrelevant at these rates rather than as unnecessary in principle.
The estimator, whose attitude is a measured input, normalises it before use
(Section~\ref{sec:frame:redmodel}).

Neither \eqref{eq:kin} nor \eqref{eq:quat} contains a mass, an inertia, or any
other physical parameter. \emph{The payload therefore enters the model only
through the dynamics.}

\subsection{Translational dynamics}\label{sec:model:trans}

Newton's second law in $\mathcal{W}$, with thrust rotated from the body frame and
a linear aerodynamic drag term, gives
\begin{equation}
\dot{\bm{v}} = \frac{1}{\mT}\Bigl( R(\bm{q})\, \ev\, T - k_d\, \bm{v} \Bigr) - g\, \ev,
\label{eq:trans}
\end{equation}
where $\mT$ is the \emph{composite} mass of airframe plus payload, $k_d$ is a
fixed drag coefficient, and $g$ is the gravitational acceleration. The mass ---
and only the mass --- enters here, as the scalar gain $1/\mT$ from thrust to
acceleration. If the thrust is known independently of the model, the measured
motion determines $\mT$ through \eqref{eq:trans} in ordinary near-hover flight,
with no dedicated excitation manoeuvre; the qualification ``independently of the
model'' is the subject of Section~\ref{sec:model:phys}.

Projecting \eqref{eq:trans} onto $\ev$ and neglecting the vertical component of
the drag term gives the vertical force balance
\begin{equation}
\mT\bigl(g + a_z\bigr) \;=\; T\cos\phi\,\cos\theta ,
\label{eq:vbal}
\end{equation}
with $\cos\phi\cos\theta$ from \eqref{eq:costheta} and $a_z=\dot v_z$. It is
used to initialise the estimator (Section~\ref{sec:model:seed}).

\paragraph{Reference point of the translational state.}\label{sec:model:refpoint}
Equation~\eqref{eq:trans} is Newton's law for the centre of mass of the
composite body, whereas $\bm p$ and $\bm v$ in \eqref{eq:state} are the position
and velocity of the body origin, the point the odometry reports. The two
coincide only for the bare airframe. With a payload attached the composite
centre of mass lies at the body-frame offset $\rc$ of \eqref{eq:rc}, and the
transport theorem for a rigid body gives
\begin{equation}
\bm v_{C} = \bm v + R(\bm q)\,\bigl(\bm\omega\times\rc\bigr),
\qquad
\dot{\bm v}_{C} = \dot{\bm v} + R(\bm q)\,\bigl(\dot{\bm\omega}\times\rc
+ \bm\omega\times(\bm\omega\times\rc)\bigr).
\label{eq:transport}
\end{equation}
Writing \eqref{eq:trans} for the origin therefore omits the second term of
$\dot{\bm v}_C$. With $\rc=c\,\ev$ and $c=(\mP/\mT)\,r_z$, its body-frame
components are
$c\,\bigl[\dot\omega_y+\omega_x\omega_z,\; -\dot\omega_x+\omega_y\omega_z,\;
-(\omega_x^2+\omega_y^2)\bigr]^{\!\top}$: the horizontal part is driven by the
angular acceleration and averages out over a manoeuvre, and the part along the
thrust axis, which is what biases a mass estimate, is quadratic in the roll and
pitch rates. Both scale with $|c|$, at most \SI{0.034}{\meter} here. On the
\SI{4}{\meter\per\second} figure-eight of Chapter~\ref{ch:evaluation} the
measured body rates and angular accelerations had root-mean-square values of
\SI{0.57}{\radian\per\second} and \SI{0.60}{\radian\per\second\squared}; the
omitted acceleration then has a root-mean-square value of
\SI{0.010}{\meter\per\second\squared}, about \SI{0.1}{\percent} of $g$, and a mean
component along the thrust axis of
\SI{-2e-4}{\meter\per\second\squared}, about \SI{0.002}{\percent} of $g$ (medians
over the $24$ flights at that speed in Section~\ref{sec:eval:reduced}). That is two
to three orders of magnitude below the manoeuvre-induced mass bias reported
there, and replaying the estimator on the composite-centre velocity of
\eqref{eq:transport} instead of the origin velocity left its mass error
unchanged (Section~\ref{sec:eval:reduced}). The model therefore keeps the origin
as its reference point and omits the transport term; it is listed with the other
approximations in Section~\ref{sec:model:approx}. The omission is specific to
this geometry and these manoeuvres: the term grows with $\mP/\mT$, with the arm
and with the square of the attitude rate, and a longer suspension or aggressive
attitude manoeuvres would have to model it.

\subsection{Rotational dynamics}\label{sec:model:rot}

Let $\Jb = \mathrm{diag}(J_{xx}, J_{yy}, J_{zz})$ denote the inertia of the bare
airframe about its own centre of mass, and $\Jt$ the composite inertia derived in
Section~\ref{sec:model:paxis}. Euler's equation in the body frame is
\begin{equation}
\dot{\bm{\omega}} = \Jt^{-1}\Bigl( \bm{\tau} - \bm{\omega} \times \Jt \bm{\omega} \Bigr),
\label{eq:rot}
\end{equation}
where $\bm{\omega} \times \Jt \bm{\omega}$ is the gyroscopic coupling.
Equation~\eqref{eq:rot} holds in this form only when the inertia and the applied
moments refer to the same point, and for a free-flying body that point must be
the centre of mass. The composite inertia must therefore be taken about the
composite centre of mass, which \emph{moves} when a payload is attached
(Section~\ref{sec:model:paxis}). For the centred payload of assumption~2 that
centre of mass stays on the thrust axis, so thrust contributes no moment to
\eqref{eq:rot}.

Taken together, \eqref{eq:trans} and \eqref{eq:rot} give each payload-dependent
quantity exactly one point of entry: the composite mass appears only as $1/\mT$
in the translational dynamics, and the inertia increment only inside $\Jt$ in
the rotational dynamics. This separation is what allows the two effects to be
attributed to distinct physical channels, and it is used throughout
Section~\ref{sec:model:onedof}.

\section{Payload-Induced Mass and Inertia Coupling}\label{sec:model:payload}

With the payload of assumption~2, define the composite mass and the
\emph{reduced mass}
\begin{equation}
\mT = \mB + \mP,
\qquad
\mu = \frac{\mB\, \mP}{\mB + \mP} = \frac{\mB\, \mP}{\mT}.
\label{eq:mu}
\end{equation}

\subsection{Inertia increment and the reduced mass}\label{sec:model:paxis}

The inertia increment requires more care than the parallel-axis theorem applied
naively, because \emph{the centre of mass itself moves when the payload is
attached}. With the airframe centre of mass at the body origin, the composite
centre of mass lies on the thrust axis at
\begin{equation}
\rc = \frac{\mP}{\mT}\, \rp = \Bigl[0,\; 0,\; \frac{\mP}{\mT}\,r_z\Bigr]^{\!\top}.
\label{eq:rc}
\end{equation}
The theorem must be applied about this \emph{new} centre of mass, and both bodies
contribute, since the airframe no longer sits at the centre of mass either.

Consider a roll or pitch axis and let $d=|r_z|$ denote the lever arm. From
\eqref{eq:rc}, the distances of the two bodies from the composite centre of mass
are
\begin{equation}
\underbrace{\bigl\| \bm{0} - \rc \bigr\| = \frac{\mP}{\mT}\, d}_{\text{airframe}},
\qquad
\underbrace{\bigl\| \rp - \rc \bigr\| = \Bigl( 1 - \frac{\mP}{\mT} \Bigr) d = \frac{\mB}{\mT}\, d}_{\text{payload}} .
\label{eq:dists}
\end{equation}
The payload's distance is \emph{smaller} than $d$. Adding the two
parallel-axis contributions and simplifying,
\begin{align}
\dJ
&= \mB \Bigl( \frac{\mP}{\mT} d \Bigr)^{\!2}
 + \mP \Bigl( \frac{\mB}{\mT} d \Bigr)^{\!2} \nonumber \\[4pt]
&= \frac{d^2}{\mT^2}\Bigl( \mB\, \mP^2 + \mP\, \mB^2 \Bigr)
 = \frac{d^2}{\mT^2}\, \mB\, \mP \bigl( \mP + \mB \bigr) \nonumber \\[4pt]
&= \frac{\mB\, \mP}{\mT}\, r_z^2
 = \mu\, r_z^2 .
\label{eq:reduced}
\end{align}
The factor $(\mB + \mP)$ produced by the sum cancels one power of $\mT$ in the
denominator, and the two contributions collapse onto the reduced mass
\eqref{eq:mu} --- the same quantity that governs the two-body problem, and for
the same reason: the inertia of a rigid pair about its centre of mass depends
only on the \emph{relative} position of the two bodies.

Equation~\eqref{eq:reduced} also identifies a natural error. Writing
$\dJ = \mP r_z^2$ --- treating the payload as orbiting a fixed body origin, and
forgetting both the airframe's own contribution and the displacement of the
centre of mass --- overestimates the increment by the factor $\mT / \mB$, which
is \SI{14.5}{\percent} at the heavier payload of this work.

The same computation carried out for the full tensor gives, for an on-axis
point mass,
\begin{equation}
\Jt \;=\; \Jb + \mu\Bigl(\lVert\rp\rVert^2 I_3 - \rp\,\rp^{\!\top}\Bigr)
\;=\; \Jb + \mathrm{diag}\bigl(\mu r_z^2,\; \mu r_z^2,\; 0\bigr).
\label{eq:Jdeployed}
\end{equation}
Three consequences are used below. The increment is \emph{isotropic} in roll and
pitch, so a single scalar $\dJ=\mu r_z^2$ describes both. It is
\emph{diagonal}: every off-diagonal product of inertia contains a horizontal arm
and vanishes. And the yaw increment is exactly zero for a point mass on the yaw
axis; the parcel's own rotational inertia about its centre is the only yaw
contribution, and Section~\ref{sec:model:approx} declares its omission.

\subsection{Magnitudes on the experimental platform}\label{sec:model:magnitudes}

The factor by which the payload slows the roll and pitch response to a given
torque is the inertia ratio
\begin{equation}
\kappa \;=\; \frac{J_{xx} + \dJ}{J_{xx}} .
\label{eq:gain}
\end{equation}
For the two payloads used in the experiments, attached at the nominal arm
$r_z=\SI{-0.27}{\meter}$ on the platform of this work
($\mB = \SI{2.0643}{\kg}$, $J_{xx} = \SI{0.0142}{\kg\meter\squared}$),
equations \eqref{eq:mu}, \eqref{eq:reduced} and \eqref{eq:gain} give
\begin{equation}
\begin{array}{lccc}
 & \mu\ [\si{\kg}] & \dJ\ [\si{\kg\meter\squared}] & \kappa \\[2pt]
\mP=\SI{0.15}{\kg} & 0.140 & 0.0102 & 1.72 \\
\mP=\SI{0.30}{\kg} & 0.262 & 0.0191 & 2.34
\end{array}
\label{eq:numbers}
\end{equation}
A \SI{0.30}{\kg} parcel on a \SI{2.06}{\kg} airframe is a \SI{15}{\percent}
change of mass but a more than \emph{two-fold} change of roll and pitch inertia.
Adapting the mass alone, in the thrust channel, accounts for the first and not
the second. The increment is governed by how far \emph{below} the airframe the
load hangs, through $r_z^2$, which is why the gripper geometry, and not only the
parcel, sets the difficulty of the task.

\section{Single-Parameter Estimation Structure}\label{sec:model:onedof}

\subsection{One online degree of freedom}

Equations \eqref{eq:mu}--\eqref{eq:Jdeployed} express every payload-dependent
quantity of the model through two numbers, the payload mass $\mP$ and the
vertical arm $r_z$, and the two are not on an equal footing.

The mass changes from flight to flight and within a flight: it is zero before a
grasp, takes the parcel's value after it, and returns to zero at release,
possibly in several steps if a load is lost piecewise. It is also
\emph{observable}, through the gain $1/\mT$ of \eqref{eq:trans}. It is therefore
the quantity estimated online, in the form of the composite mass $\mT$, and the
payload mass follows as
\begin{equation}
\mP = (\mT - \mB)^{+},
\label{eq:mPplus}
\end{equation}
where $(\cdot)^{+}$ denotes the positive part, because the mass estimate is
permitted to fall slightly below $\mB$ to leave room for calibration error.

The vertical arm, by contrast, is a property of the gripper rather than of the
parcel. It is fixed by where the gripper holds the load relative to the rotor
plane, and it enters the equations only through the inertia increment, which is
weakly excited in delivery flight. It is therefore held at a nominal value
rather than estimated. The inertia increment consumed by the model is then a
function of the single online scalar,
\begin{equation}
\boxed{\;\dJ(\mT) \;=\; \frac{\mB\,(\mT-\mB)^{+}}{\mT}\;r_z^2\;}
\label{eq:dJofm}
\end{equation}
which vanishes when the mass estimate returns to the bare-airframe value. No
attach or release event is needed to switch the inertia on or off: the algebra
does it, continuously, as the estimate moves. The mass is estimated; the inertia
is derived from it and is never an independent decision variable.

\subsection{Which channel identifies the mass}\label{sec:model:rotchannel}

Equation~\eqref{eq:dJofm} makes the rotational dynamics depend on $\mT$. Should
the \emph{estimator} also be allowed to use that dependence, so that the
angular-rate residual contributes information about the mass?

In principle it should: the increment is larger than the airframe inertia
itself (Section~\ref{sec:model:magnitudes}), so any roll or pitch manoeuvre
carries a strong mass signature. In practice the rotational channel of this
platform is not modelled well enough to carry an absolute quantity. Three
independent errors were found in it: the calibrated airframe inertia differs
substantially from the simulator's model description; rotor drag acting
through the vertical offset of the rotor plane produces a velocity-dependent body
moment that \eqref{eq:rot} does not contain; and the reconstructed torque of an
unloaded vehicle carries a small constant bias
(Section~\ref{sec:model:approx}). None prevents the vehicle from flying, but an
optimiser that is allowed to move $\mT$ inside the rotational model will use the
mass to absorb all three. In development flights with that freedom the estimator
overestimated the loaded mass by \SIrange{28}{64}{\percent} at
\SI{4}{\meter\per\second} and the flight diverged; with the airframe inertia
corrected to the simulator's value it still overestimated the mass by
\SI{29}{\percent} in the figure-eight and by \SI{39}{\percent} after release, so
that the release was never confirmed.

The estimator therefore does not integrate the rotational dynamics at
all. It keeps the translational dynamics \eqref{eq:trans}, supplies the measured
attitude as a known input instead of estimating it, and holds the composite mass
and the velocity as its only states (Section~\ref{sec:frame:redmodel}). The mass
then enters the estimator's model only where the model --- a thrust, a mass and
gravity --- is trustworthy, and no inertia, torque or arm enters it.

Removing only the \emph{direct} sensitivity of the rotational dynamics to the
mass would not be enough. An estimator that integrates the rotational dynamics
keeps two indirect couplings. Within a window, the angular-rate residual shapes
the estimated attitude, and the attitude sets the direction of the thrust in
\eqref{eq:trans}, so an error in the rotational model reaches the mass through
the attitude. And any payload inertia it carries in its rotational model must
come from a mass estimate, so a biased estimate changes the next window's
rotational model. The estimator of this work has neither path: it has no
estimated attitude through which an error could pass, and no inertia into which a
mass could be fed back. Section~\ref{sec:eval:reduced} compares it with a
full-state estimator that removed only the direct path
(Section~\ref{sec:model:mhe}).

\subsection{What is estimated, what is derived, and what is fixed}\label{sec:model:whynot}

Table~\ref{tab:params} records the identity of every parameter of the model. One
scalar is an online degree of freedom; everything else is derived from it by the
algebra of Section~\ref{sec:model:payload}, calibrated once for the airframe, or
held at a nominal value.

\begin{table}[ht]
\caption{Model parameters, classified by identity. Only the composite mass is a
decision variable of the estimator; the inertia increment is generated from it.}
\label{tab:params}
\centering
\footnotesize
\begin{tabular}{@{}l p{0.20\textwidth} c p{0.26\textwidth} p{0.16\textwidth}@{}}
\toprule
Symbol & Meaning & Dim. & Identity & Value / enters \\
\midrule
$\mT$ & composite mass & 1 & \textbf{MHE decision variable} & \eqref{eq:trans} \\
\midrule
$\dJ$ & roll/pitch inertia increment & 1 & derived, \eqref{eq:dJofm}; absent from the estimator & \eqref{eq:predmodel} \\
$\mP$ & payload mass & 1 & derived, $(\mT-\mB)^{+}$ & \eqref{eq:dJofm} \\
\midrule
$r_z$ & payload vertical arm & 1 & nominal gripper geometry & \SI{-0.27}{\meter} \\
$\cvec$ & CoM horizontal offset & 2 & \textbf{not modelled}; held at $\bm 0$ & --- \\
\midrule
$\mB$ & bare-airframe mass & 1 & calibration & \SI{2.0643}{\kg} \\
$\Jb$ & bare-airframe inertia & 3 & calibration & $\mathrm{diag}(0.0142,$ $0.0142, 0.0210)$ \si{\kg\meter\squared} \\
$k_f, k_m$ & thrust and yaw-reaction coefficients & 2 & calibration & \SI{8.55e-6}{\newton\second\squared}, \SI{0.016}{\meter} \\
$k_d$, $g$ & drag coefficient, gravity & 2 & constants & \SI{0.05}{\kg\per\second}, \SI{9.81}{\meter\per\second\squared} \\
\bottomrule
\end{tabular}
\end{table}

Three entries of the table invite objection.

\paragraph{Why not augment the inertia increment as well?} Because it is nearly
unobservable in the flight regime of interest, and because the channel through
which it would be observed is the one this platform models least well. The
increment acts only through $\dot{\bm{\omega}}$, which a delivery task spent
largely in near-hover, low-rate flight excites weakly; a free $\dJ$ would let the
optimiser explain measurement noise and the rotational model errors of
Section~\ref{sec:model:rotchannel} with a quantity the data barely constrain.
Moreover $\Jt$ appears \emph{inverted} in \eqref{eq:rot} and is propagated across
the horizon, so an error in it is amplified rather than damped, and inertia
elements are two orders of magnitude smaller than the mass, which would worsen
the scaling of the estimator's Hessian.

\paragraph{Why not estimate a horizontal offset?} Because at the operating
condition of this work there is none to estimate, and carrying one is costly. The
gripper aligns over the parcel before attaching, so the offset is bounded by the
alignment tolerance; an estimated offset would be identifiable only through the
rotational residual, the channel that Section~\ref{sec:model:rotchannel} keeps
away from absolute quantities; and an offset is undefined once the payload has
gone, so any parameterisation that carries one must also say what it becomes at
release.

\paragraph{Why is a nominal $r_z$ acceptable?} Because the arm is a property of
the gripper, fixed by its design, and not information about the parcel. It is a
geometric prior, and the method does use it: it enters the inertia increment
\eqref{eq:dJofm} of the controller, where its error propagates quadratically ---
a \SI{10}{\percent} error in $r_z$ is a \SI{21}{\percent} error in $\dJ$, which
mis-sizes the predicted attitude response and the inner-loop ratio $\kappa$. It
does not enter the estimator at all. The difference between the nominal and the realised arm is quantified
in Section~\ref{sec:model:approx}.

\paragraph{The information boundary.} The precise claim of this dissertation is
that \emph{no task-specific payload mass enters the estimator or the controller,
and no attach or release notification enters the estimator}. It is not the
stronger claim that the flight stack contains no payload-related constant. Two
such constants remain and are disclosed: the nominal arm $r_z$ above, and the
platform's \SI{0.30}{\kg} payload \emph{envelope}, which sets the inertia floor
of the attach-time bootstrap (Section~\ref{sec:frame:bootstrap}). Both are
specifications of the airframe and its gripper, in the same category as $\Jb$ or
the actuator envelope, and neither is the mass of the parcel carried on a
particular flight. The controller's own release command is used to open the
release confirmation, to start its timeout, to end the inertia floor and to
classify an acceptance as a delivered rather than a lost load
(Chapter~\ref{ch:framework}); it does not set the payload mass.

\section{Measured-Input Reconstruction and Model Interfaces}\label{sec:model:phys}

\subsection{The physical input}

The thrust and torque in \eqref{eq:trans} and \eqref{eq:rot} must be known to the
estimator. They are reconstructed from measured rotor speeds $\omega_i$ through
the quadratic rotor law $F_i = k_f \omega_i^2$ and the linear map from the four
rotor thrusts to the wrench. For the X configuration of this platform, with
rotor positions $(x_i,y_i)$, $|x_i|=|y_i|=l=\SI{0.174}{\meter}$, and spin
directions $d_i\in\{\pm1\}$,
\begin{equation}
\begin{bmatrix} \Tphys \\ \tauphys \end{bmatrix}
=
\underbrace{
\begin{bmatrix}
1 & 1 & 1 & 1 \\
-l & \phantom{-}l & \phantom{-}l & -l \\
-l & \phantom{-}l & -l & \phantom{-}l \\
-k_m & -k_m & \phantom{-}k_m & \phantom{-}k_m
\end{bmatrix}}_{M}
\begin{bmatrix} F_1 \\ F_2 \\ F_3 \\ F_4 \end{bmatrix},
\qquad F_i = k_f \omega_i^2 ,
\label{eq:phys}
\end{equation}
whose rows are $\Tphys=\sum_i F_i$, $\tau_x=\sum_i y_iF_i$,
$\tau_y=-\sum_i x_iF_i$ and $\tau_z=-k_m\sum_i d_iF_i$, with
$k_f=\SI{8.55e-6}{\newton\second\squared}$ and yaw-reaction coefficient
$k_m=\SI{0.016}{\meter}$. The first two moment components are lever-arm
products; the third is the aerodynamic reaction torque, proportional to rotor
thrust and opposed to the rotor's own spin.

\paragraph{This is a wrench, not a state estimate.} Equation~\eqref{eq:phys}
returns the wrench the \emph{actuators generate}. It excludes $\Jt\dot{\bm\omega}$
and rotor gyroscopic terms, which is what makes it usable as the known input
$u^{\mathrm{phys}}$ of the estimator rather than as a second, competing estimate
of the state. The estimator uses only its first row, $\Tphys$. Rotor speeds are
sampled at \SI{250}{\hertz} and averaged over each \SI{10}{\hertz} estimation
interval, consistent with a zero-order-hold discretisation of \eqref{eq:trans}.

\paragraph{Why the commanded input cannot be used.} Let the controller compute
its thrust command from an assumed mass $\tilde{m}$, so that in hover
$T_{\mathrm{cmd}} = \tilde{m} g$. Feeding $T_{\mathrm{cmd}}$ into
\eqref{eq:trans} and solving for the mass yields
$\mT = T_{\mathrm{cmd}} / g = \tilde{m}$ for \emph{any} $\tilde{m}$: the relation
is satisfied identically, the residual carries no information, and the estimate
would reproduce whatever the controller already believed. Any inversion error,
saturation or motor lag between command and rotor is likewise invisible.
Equation~\eqref{eq:phys} breaks this circularity because rotor speeds are a
physical measurement into which no model state enters, taken downstream of
allocation and the motor dynamics. The principle recurs throughout the design: an
adaptive loop must be fed by a signal that is not itself a product of the
adaptation. The difference is not academic; the executed yaw torque recorded in
flight is about a third of the intended value (Section~\ref{sec:eval:platform}).

\subsection{The interface to the autopilot}\label{sec:model:interface}

The deployed system is a cascade: the predictive controller is the outer loop,
and the autopilot's body-rate controller and mixer are the inner loop. Three
torque-like quantities exist in it and must not be confused.
\begin{enumerate}
\item The \emph{predicted torque} $\bm\tau$ is a decision variable of the
controller's problem. It shapes the predicted state trajectory through
\eqref{eq:predmodel} and never leaves the solver.
\item The \emph{body-rate setpoint} is what the autopilot receives. After each
solve the controller publishes, at \SI{50}{\hertz}, a collective thrust and a
body-rate setpoint interpolated along the latest predicted trajectory: the
thrust from the predicted input sequence, converted to a normalised throttle by
\eqref{eq:throttleinv}, and the body rate from the predicted state sequence one
stage ahead, starting at its second node. The roll and pitch rate setpoints are
then scaled for the payload inertia (Section~\ref{sec:frame:omegascale}). The
autopilot's rate loop and mixer turn this setpoint into rotor commands.
\item The \emph{physical torque} $\tauphys$ is what the rotors actually
produced, reconstructed from measured rotor speeds by \eqref{eq:phys}. No
optimiser uses it; it is logged as a diagnostic.
\end{enumerate}
The division is deliberate: the rate loop runs at a bandwidth the
\SI{20}{\hertz} solver cannot match, and leaving allocation on the autopilot
keeps saturation handling next to the hardware. It is also why the estimator's
thrust is evaluated in the \emph{forward} direction from measured rotor speeds
rather than read off the solver: $M$ applied to what the motors did is a
measurement, whereas the solver's input is an intention. The ROS layer
never evaluates the inverse allocation; its role in sizing the controller's
input bounds is discussed with the actuator envelope in
Section~\ref{sec:frame:inputs}.

Figure~\ref{fig:paramflow} traces the complete accounting: what is measured,
what is estimated, what is generated by algebra, where each quantity is consumed,
and how the loop closes.

\begin{figure}[tbp]
\centering
\begin{tikzpicture}[
  x=1cm,y=1cm,font=\scriptsize,
  nd/.style={draw=gray!70,rounded corners=1.5pt,align=center,inner sep=3pt,
             font=\scriptsize},
  key/.style={nd,draw=black,line width=0.7pt,fill=gray!15},
  der/.style={nd,fill=gray!4},
  gh/.style={nd,draw=gray!50,text=gray!40!black},
  ar/.style={-{Latex[length=1.5mm]},semithick,gray!25!black},
  ln/.style={semithick,gray!25!black},
  bus/.style={-{Latex[length=2mm]},line width=0.9pt,gray!15!black},
  dar/.style={-{Latex[length=1.4mm]},dashed,thin,gray!55!black},
  lb/.style={font=\tiny,align=center,inner sep=1.2pt},
  lane/.style={draw=gray!45,dashed,rounded corners=2.5pt,inner sep=3mm},
  hd/.style={font=\tiny\bfseries,text=gray!50!black,inner sep=1.2pt},
]

%========================= band 1: measured =========================
\node[nd]  (rot) at (2.0,0)   {rotor speeds $\omega_i$\\ \SI{250}{\hertz}};
\node[der] (wr)  at (5.8,0)   {$\Tphys,\ \tauphys$\\ \eqref{eq:phys}};
\node[nd]  (odo) at (10.6,0)  {odometry\\ $\bm p,\bm v,\bm q,\bm\omega$};
\draw[ar] (rot) -- node[lb,above]{$M$ forward} (wr);
\begin{scope}[on background layer]
  \node[lane,fit=(rot)(wr)(odo)] (L1) {};
\end{scope}
\node[hd,anchor=south west] at ([xshift=1mm]L1.north west) {MEASURED --- downstream of allocation, independent of every model state};

%========================= band 2: estimated ========================
\node[nd,text width=62mm] (res) at (6.6,-2.1)
  {model \eqref{eq:redmodel} with measured attitude $\tilde{\bm q}$;\ residual $\bm e^v_i=\bm y^v_i-\bm v_i\in\mathbb{R}^{3}$ \eqref{eq:redresid}};
\node[nd]  (vres) at (4.8,-3.35) {velocity residual\\ $\bm e^v_i$};
\node[key] (mT)   at (4.8,-4.5)  {$\mT$};
\node[gh]  (seed) at (1.5,-4.5)  {seed\\ \eqref{eq:winbal}};
\node[gh,text width=26mm]  (prior)at (12.7,-3.9) {fixed inputs\\ $r_z$ nominal\\ $\mB,\Jb$ calibrated};
\draw[ar] (wr.south) -- (res.north -| wr.south);
\node[lb,anchor=east] at (5.65,-0.95) {$\Tphys$ (input)};
\draw[ar] (odo.south) -- ++(0,-1.05) -| (res.north -| 8.6,0);
\node[lb,anchor=west] at (10.75,-0.95) {$\bm y_i$ (measured)};
\draw[ar] (res.south -| vres.north) -- (vres.north);
\draw[ar] (vres) -- (mT);
\draw[dar] (seed) -- node[lb,above]{initial\\ iterate only} (mT);
\begin{scope}[on background layer]
  \node[lane,fit=(res)(vres)(mT)(seed)(prior)] (L2) {};
\end{scope}
\node[hd,anchor=south west] at ([xshift=1mm]L2.north west) {ESTIMATED --- the only online degree of freedom};

\draw[bus] (mT.south) -- node[lb,right,pos=0.55]{one scalar, nothing else} (4.8,-6.4);

%========================= band 3: derived ==========================
\node[der] (mp)  at (3.0,-7.4)  {$\mP=(\mT-\mB)^{+}$\\ \eqref{eq:mPplus}};
\node[der] (dj)  at (7.6,-7.4)  {$\dJ=\mu(\mT)\,r_z^2$\\ \eqref{eq:dJofm}};
\node[gh]  (cz)  at (11.9,-7.4) {$\cvec\equiv\bm 0$\\ centred grasp};
\draw[ar] (mp) -- (dj);
\begin{scope}[on background layer]
  \node[lane,fit=(mp)(dj)(cz)] (L3) {};
\end{scope}
\node[hd,anchor=south west] at ([xshift=1mm]L3.north west) {DERIVED --- rigid-body algebra, no new unknowns};
\draw[dar] (prior.south) -- ++(0,-1.5) -| node[lb,pos=0.22,above]{$\to \dJ$} (dj.north);
\draw[ar] (1.0,-7.4) -- (mp.west);

\draw[ln] (4.8,-6.4) -- (1.0,-6.4) -- (1.0,-8.6) -- (4.4,-8.6);
\fill[gray!25!black] (1.0,-7.4) circle (0.6pt);
\node[lb,rotate=90,anchor=south] at (0.92,-7.4) {$\mT$ passes straight through};
\draw[ar] (dj.south)  -- (7.6,-8.6);
\draw[ln] (4.4,-8.6) -- (7.6,-8.6);
\draw[bus] (6.1,-8.6) -- node[lb,right,pos=0.6]
  {$p_{\mathrm{NMPC}}=[\,x_{\mathrm{ref}};\mT;\dJ;\bm 0;\hat{\bm d}\,]\in\mathbb{R}^{20}$ \eqref{eq:pnmpc}}
  (6.1,-9.6);

%========================= band 4: consumed =========================
\node[nd,text width=54mm,align=left] (mdl) at (4.4,-10.7)
  {\textbf{prediction model} $F$ \eqref{eq:predmodel}\\[1pt]
   $1/\mT \to \dot{\bm v}$\\
   $\Jt(\dJ) \to \dot{\bm\omega}$, diagonal};
\node[nd,text width=44mm,align=left] (cst) at (11.0,-10.7)
  {\textbf{cost baseline} $u_{\mathrm{hover}}$ \eqref{eq:hover}\\[1pt]
   $T_h=\mT(g-\hat d_z)$\\ torque entries $=0$};
\begin{scope}[on background layer]
  \node[lane,fit=(mdl)(cst)] (L4) {};
\end{scope}
\node[hd,anchor=south west] at ([xshift=1mm]L4.north west) {CONSUMED --- the same numbers enter the constraint and the objective};

%========================= band 5: actuated =========================
\node[nd] (eta) at (5.0,-12.7) {throttle $\eta(T)$ \eqref{eq:throttleinv}};
\node[nd] (w1)  at (5.0,-13.6) {body-rate $\bm\omega_1$, scaled by $\kappa(\dJ)$};
\node[nd] (px4) at (11.0,-13.15) {PX4 rate loop $+$ mixer \eqref{eq:mixerinv}};
\draw[ar] (eta.east) -- (px4.west|-eta.east) -- (px4.west);
\draw[ar] (w1.east)  -- (px4.west|-w1.east)  -- (px4.west);
\begin{scope}[on background layer]
  \node[lane,fit=(eta)(w1)(px4)] (L5) {};
\end{scope}
\node[hd,anchor=south east] at ([xshift=-1mm]L5.north east) {ACTUATED --- $\bm\tau$ never leaves the solver};
\draw[ln] (mdl.south) -- ++(0,-0.55) -| (2.0,-12.7);
\draw[ar] (2.0,-12.7) -- (eta.west);
\draw[ln] (2.0,-12.7) -- (2.0,-13.6);
\draw[ar] (2.0,-13.6) -- (w1.west);

%========================= loop closure =============================
\draw[bus,gray!40!black] (px4.east) -- (14.4,-13.15) -- (14.4,1.55)
  -- node[lb,above,pos=0.55]{the loop closes through \emph{measurement}, never through the controller's intention} (2.0,1.55) -- (rot.north);
\end{tikzpicture}
\caption[How the payload parameter reaches the controller]{How the payload parameter reaches the controller. In the
estimator the window's measured velocities and its known inputs --- the
measured attitude and the reconstructed thrust --- meet in one velocity residual
\eqref{eq:redresid}, which identifies the single boxed decision variable. Everything in the \textsc{derived} band follows
from $\mT$ by the algebra of Section~\ref{sec:model:payload} and introduces
neither a new unknown nor a new sensor; the inertia increment vanishes by itself
when the estimate returns to the bare airframe, so no release event is needed to
remove it. The estimator has no rotational model. The horizontal-offset parameters are held at zero by the operating
condition, not estimated. The loop closes through measured rotor speeds, never
through the moment the solver intended. The residual specific force $\hat{\bm d}$ in the controller's parameter vector is estimated from the same measurements but is not a payload parameter (Section~\ref{sec:model:residual}).}
\label{fig:paramflow}
\end{figure}

\section{Modelling Assumptions and Limitations}\label{sec:model:approx}

The relations of this chapter fall into two classes, and the distinction matters
for how the results should be read.

\paragraph{Relations that hold exactly.} For a payload on the thrust axis, the
vertical displacement \eqref{eq:rc} of the centre of mass is collinear with the
thrust, so the thrust produces no moment about the composite centre of mass;
this is a consequence of geometry, not an omitted term. The increment
\eqref{eq:reduced} is exact for a rigid pair of point masses, and so is the zero
yaw increment noted after \eqref{eq:Jdeployed}.

The remaining relations are approximations whose validity depends on the
operating condition.

\paragraph{The payload is assumed centred.} The model sets the horizontal
attachment offset to zero. A residual offset $r_{xy}$ left by imperfect
alignment produces a constant trim moment $\mP g\, r_{xy}$ and a roll/pitch
inertia error $\mu r_{xy}^2$. At the \SI{0.03}{\meter} attachment tolerance of
the gripper and \SI{0.30}{\kg} these are \SI{0.09}{\newton\meter}, or
\SI{18}{\percent} of the roll authority, and \SI{0.4}{\percent} of the
increment. The first is supplied by the integral action of the autopilot's rate
loop; the second is negligible.

\paragraph{The arm is nominal, not measured.} Both optimisers use
$r_z=\SI{-0.27}{\meter}$ for every flight. The realised arm is set by the height
at which the vehicle grasps the parcel, which is commanded in the autopilot's
local frame whose vertical zero drifts from flight to flight; it varies between
about \SI{0.25}{\meter} and \SI{0.30}{\meter}. Over that range the true increment
for \SI{0.30}{\kg} lies between \SI{0.0164}{} and
\SI{0.0236}{\kg\meter\squared} against the modelled
\SI{0.0191}{\kg\meter\squared} --- an error of $-14$ to $+23$ percent, and a
true inertia ratio between $2.15$ and $2.66$ against the modelled $2.34$. Every
result in this dissertation carries an arm error of this size, unmeasured on a
per-flight basis.

\paragraph{The translational state refers to the body origin.} Equation
\eqref{eq:trans} is applied to the body origin rather than to the composite
centre of mass, which omits the transport acceleration of \eqref{eq:transport}.
Its thrust-axis component is about \SI{0.002}{\percent} of $g$ in the
\SI{4}{\meter\per\second} figure-eight, and removing the omission in replay did
not change the mass error (Section~\ref{sec:model:refpoint}).

\paragraph{The parcel's own inertia is neglected.} The payload is treated as a
point mass. A box of side \SI{0.15}{\meter} has a central inertia of about
$0.004\,\mP$, about \SI{2}{\percent} of the roll/pitch increment and the only
yaw contribution of a centred load.

\paragraph{The airframe rotational model is approximate.} The calibrated
airframe inertia $J_{xx}=\SI{0.0142}{\kg\meter\squared}$ is about
\SI{60}{\percent} of the value in the simulator's model description, and a
least-squares fit of bare-airframe pitch dynamics recovered an effective inertia
$1.07$--$1.18$ times the latter; and the rotor-drag moment and the constant
torque bias of order \SI{0.01}{\newton\meter} described in
Section~\ref{sec:model:rotchannel} are not modelled. These errors are the
reason the mass is kept off the rotational channel, and the reason the
estimator does not integrate that channel at all
(Section~\ref{sec:model:rotchannel}). They remain in the controller's prediction
model, where they appear as an unmodelled moment that the rate loop absorbs.

\paragraph{The inertia is treated as constant across a transfer.}
Equation~\eqref{eq:rot} omits the $\dot{\Jt}\bm{\omega}$ term that a
time-varying inertia would contribute. The omission is exact on either side of a
grasp or release and wrong only for the duration of the transfer itself, which
the controller's prediction shares with the plant's rate loop.

\paragraph{The rotor model is static.} The thrust coefficient $k_f$ is a
constant, and is known to lose validity under large induced inflow during fast
translation \cite{bangura2012}; the mass channel recovers $\mT$ from an absolute
thrust and inherits any error in $k_f$ in full. The yaw reaction of
\eqref{eq:phys} is the steady-state term only, omitting the rotor
angular-acceleration contribution $\sum_i J_r \dot\omega_i$, and rotor
gyroscopic moments are not modelled. These terms are negligible in hover and
mild manoeuvre but not during rapid yaw. Because $\mT$ is recovered from an
absolute thrust, a relative error in $k_f$ passes one-for-one into $\hat m$; this
sensitivity is analytical and has not been measured in the evaluated
configuration (Section~\ref{sec:eval:validity}). Translational data cannot
separate the two: the estimator flags a thrust-map error that drives
its estimate outside a physical band (Section~\ref{sec:frame:health}) but does
not correct it.

\section{Chapter Summary}\label{sec:model:summary}

This chapter addressed the first research objective and the structural part of
the second. For a rigid payload grasped on the thrust axis, the roll and pitch
inertia increment is the reduced mass times the square of the vertical arm, so
the composite mass is the only payload quantity that has to be identified
online; the inertia follows from it algebraically and vanishes with it. The mass
is identified on the translational channel, from a thrust reconstructed from
measured rotor speeds. The estimator does not integrate the rotational
dynamics at all --- the measured attitude enters as a known input --- so the
optimiser cannot use the mass to absorb rotational model error. The structure rests on stated
assumptions --- a centred grasp, a nominal arm and a calibrated rotor model ---
whose limits are collected in Section~\ref{sec:model:approx}. What remains open
is the numerical form of the estimator and the controller, and how a release is
recognised; Chapter~\ref{ch:framework} takes these up.

\end{spacing}
\newpage
